\chapter{Instantiating Biophilic Interaction Through AI-Transformed Sound}
\label{ch:moss}

\begin{synopsis}
\synopsislead{What happens when the soundscape of an office becomes}{the material for biophilic interaction? The preceding chapters expanded and structured the possibilities of multisensory biophilic design. This chapter develops one of those possibilities as a concrete design instance, using sound to examine how an interactive system might alter the affective character of an indoor environment while retaining a relationship with the activity already taking place within it.}

This chapter introduces the Mediated Office Soundscapes System (MOSS), a headphone-based research prototype that uses neural audio transformation to render ongoing office sound through bird-like or water-like auditory textures. Rather than introducing an independent recording of nature, MOSS uses the acoustic activity of the office itself as source material. It is designed around three experiential functions: enriching quietness, reframing office activity, and preserving environmental connection.

The chapter examines this proposition through two complementary studies. An online listening study compares Office-Only, Conventional-Nature, and AI-Transformed soundscapes across calm and active office scenes, examining pleasantness, calmness, nature-likeness, distraction, work compatibility, contextual fit, and participants' interpretations of the transformed sound. A subsequent live prototype study with professionals in architecture, workplace design, and spatial sound examines MOSS as an interactive and spatial intervention, including questions of responsiveness, intelligibility, control, agency, environmental awareness, and architectural integration.

Within the trajectory of the dissertation, this chapter addresses \rqref{sound} by bringing the affective and biophilic directions developed across the preceding chapters into a single design instance. MOSS treats environmental sound as something that can be interpreted and transformed rather than simply reduced, masked, or replaced. In doing so, the chapter examines how AI-enabled mediation might support a different relationship between occupants and a shared indoor soundscape, and what experiential and design conditions are required for such mediation to become part of everyday workplace environments.
\end{synopsis}

\chaptersource{%
This chapter is based on: Shruti Rao, Judith Good, and Hamed Alavi.
\emph{Exploring the Potential of an AI-Transformed System for Biophilic Soundscapes in Future Office Spaces}.
Manuscript in preparation.%
}

\clearpage

% \import{papers/05-biophilic-sound/}{thesis-body.tex}
\clearpage

\begin{chapterclosing}{Concluding synthesis}
This chapter develops AI-transformed sound as a concrete instance of affective biophilic interaction. MOSS uses ongoing office activity as material for a nature-like auditory layer, allowing environmental events to continue shaping what occupants hear while altering the experiential character through which those events are encountered. The resulting design proposition brings together three functions: enriching acoustically sparse periods, reframing more noticeable activity, and maintaining a perceptible relation with the surrounding office.

The listening study provides evidence that this form of mediation produces an experiential profile distinct from both the unmodified office and conventional nature playback. The AI-transformed condition was rated highest overall for calmness and lowest for distraction and cold or clinical character, while conventional nature remained more recognisably nature-like. It was also comparatively stable across calm and active office scenes. These findings indicate that the value of biophilic transformation need not depend on reproducing nature literally. Its design potential also lies in how transformation changes the affective character of existing environmental activity while retaining a relation to its source.

At dissertation level, MOSS demonstrates how affective inquiry, multisensory biophilic design, and interactive technology can meet within a concrete indoor system. It also makes visible the conditions such systems must negotiate: peripheral attention, naturalness and ambiguity, responsiveness and intelligibility, environmental awareness, occupant control, shared governance, and architectural fit. The design proposition is therefore less one of replacing an undesirable environment with a preferred one than of mediating how occupants perceive, interpret, and remain connected to the environment they already inhabit.

\chaptercarryforward{The concluding chapter returns to the dissertation as a whole, drawing together what these studies establish about affective HBI as an orientation for investigating indoor experience, developing experiential possibilities for future environments, and carrying those possibilities into concrete forms of interaction.}
\end{chapterclosing}